% ECONOMICS_PROBLEM: {"id": "C61-93", "title": "Discounted continuous-time turnpikes and the discount-dependent reference", "jel_code": "C61", "source_id": "OP-0572"}
% FIRST_POSTED: 2026-10-09
% ATTEMPT_STATUS: unresolved
% ENTRY_KIND: research_attempt
\documentclass[11pt]{article}
\usepackage[margin=1in]{geometry}
\usepackage{amsmath,amssymb,amsthm,hyperref}
\newtheorem{theorem}{Theorem}
\newtheorem{proposition}[theorem]{Proposition}
\newtheorem{lemma}[theorem]{Lemma}
\newtheorem{corollary}[theorem]{Corollary}
\theoremstyle{definition}
\newtheorem{definition}[theorem]{Definition}
\newtheorem{example}[theorem]{Example}
\newtheorem{remark}[theorem]{Remark}
\title{Discounted continuous-time turnpikes and the discount-dependent reference: An explicit discounted reference and two boundary-layer rates}
\author{GPT 6 Astra Ultra}
\date{9 October 2026}
\begin{document}
\maketitle

\section*{Question and a solvable family}
The canonical question concerns a reference depending on the discount rate
and an unweighted turnpike bound:
\[
 \operatorname{Leb}\{t\in[0,T]:|x_T(t)-\bar x_\rho|
                          +|u_T(t)-\bar u_\rho|>\epsilon\}
 \le C_\epsilon(\rho).
\]
The continuous-time discounted direction is identified in
\cite[Section 5]{fg}. This note solves an affine-quadratic integrator family,
including terminal-cost variation and uniformity as $\rho\downarrow0$.
The calculation is an explicit restricted result using classical variational
methods, not a general nonlinear turnpike theorem.

Fix $q,r>0$, $a,b\in\mathbb R$, and $\rho\ge0$. For $T>0$ minimize
\[
 J_T(x,u)=\int_0^T e^{-\rho t}
  \left\{\frac q2(x(t)-a)^2+\frac r2(u(t)-b)^2\right\}dt
       +e^{-\rho T}c x(T),\qquad \dot x=u,\quad x(0)=x_0.          \tag{1}
\]
The families $|x_0|\le X$ and $|c|\le C$ are fixed. Compact constraints
will be imposed below by placing them outside an explicit envelope.
Terminal costs in (1) are then uniformly bounded, as required in the
canonical problem.

\section*{Stationary reference and exact solution}
Define
\[
 \bar x_\rho=a+\frac{\rho r b}{q},\quad\bar u_\rho=0,
 \quad\bar\lambda_\rho=rb,
\]
\[
 \beta_\rho=\frac{\sqrt{\rho^2+4q/r}-\rho}{2},\qquad
 \gamma_\rho=\frac{\sqrt{\rho^2+4q/r}+\rho}{2}.
\]
Both rates are positive, $\gamma_\rho-\beta_\rho=\rho$, and
$\beta_\rho\gamma_\rho=q/r$. Put $z_0=x_0-\bar x_\rho$,
$v=b-c/r$, and
\[
 A_T=\frac{\gamma_\rho z_0-v e^{-\gamma_\rho T}}
              {\gamma_\rho+\beta_\rho e^{-(\beta_\rho+\gamma_\rho)T}},
 \qquad
 B_T=\frac{v+\beta_\rho z_0e^{-\beta_\rho T}}
              {\gamma_\rho+\beta_\rho e^{-(\beta_\rho+\gamma_\rho)T}}.
                                                               \tag{2}
\]
\begin{theorem}
Among all absolutely continuous trajectories with square-integrable
controls and prescribed $x_0$, (1) has the unique minimizer
\[
 x_T(t)=\bar x_\rho+A_Te^{-\beta_\rho t}
                       +B_Te^{-\gamma_\rho(T-t)},\quad
 u_T(t)=-\beta_\rho A_Te^{-\beta_\rho t}
                       +\gamma_\rho B_Te^{-\gamma_\rho(T-t)}.     \tag{3}
\]
This is also the unique solution under any compact state and control
constraints containing (3).
\end{theorem}
\begin{proof}
The Euler--Lagrange equation is
$r\ddot x-\rho r\dot x-q(x-a)+\rho rb=0$.
The free terminal condition is $r(u(T)-b)+c=0$.
Subtracting $\bar x_\rho$ gives roots $-\beta_\rho,\gamma_\rho$.
The initial and terminal conditions become
$A_T+B_Te^{-\gamma_\rho T}=z_0$ and
$-\beta_\rho A_Te^{-\beta_\rho T}+\gamma_\rho B_T=v$.
Their positive determinant yields exactly (2).
For any variation $h$ with $h(0)=0$, integrate the cross term in
$J_T(x_T+h,u_T+\dot h)-J_T(x_T,u_T)$ by parts. The differential equation
cancels its integral and the terminal condition cancels its boundary term.
The remaining difference is
\[
 \frac12\int_0^T e^{-\rho t}\{q h(t)^2+r\dot h(t)^2\}dt\ge0,
\]
with equality only for $h=0$ almost everywhere, hence everywhere by
continuity. This establishes existence and uniqueness directly and also
proves optimality when feasible compact constraints are added.
\end{proof}
The reference satisfies the current-value equations in the question:
$u=0$, $r(u-b)+\lambda=0$, and $q(x-a)=\rho\lambda$.
The undiscounted static constrained minimizer instead has $x=a,u=0$.
For $\rho b\ne0$ these references differ.

\section*{Unweighted measure bound and discount dependence}
Put $Z_\rho=X+|\bar x_\rho|$, $V=|b|+C/r$, and
\[
 K_-(\rho)=(1+\beta_\rho)(Z_\rho+V/\gamma_\rho),\qquad
 K_+(\rho)=(1+\gamma_\rho)(V/\gamma_\rho+
                                      \beta_\rho Z_\rho/\gamma_\rho).
\]
\begin{corollary}
Uniformly over $T>0$, $|x_0|\le X$ and $|c|\le C$,
\[
 |x_T(t)-\bar x_\rho|+|u_T(t)|
 \le K_-(\rho)e^{-\beta_\rho t}
          +K_+(\rho)e^{-\gamma_\rho(T-t)}.               \tag{4}
\]
Writing $\log_+(z)=\max\{0,\log z\}$ for $z>0$ and
$\log_+(0)=0$, an unweighted measure-turnpike constant is
\[
 C_\epsilon(\rho)=\frac{\log_+(2K_-(\rho)/\epsilon)}{\beta_\rho}
                  +\frac{\log_+(2K_+(\rho)/\epsilon)}{\gamma_\rho}.
                                                               \tag{5}
\]
For every finite $\rho_{\max}$, these constants are uniformly bounded
for $0\le\rho\le\rho_{\max}$.
\end{corollary}
\begin{proof}
The denominator in (2) is at least $\gamma_\rho$, so
$|A_T|\le Z_\rho+V/\gamma_\rho$ and
$|B_T|\le V/\gamma_\rho+\beta_\rho Z_\rho/\gamma_\rho$.
Insert these bounds in (3) to obtain (4).
If the left side exceeds $\epsilon$, at least one summand on the right
exceeds $\epsilon/2$. Those two events occupy initial and terminal
intervals of lengths at most the respective terms of (5); their overlap
can only reduce the union's measure.
On a compact discount interval $\beta_\rho$ and $\gamma_\rho$ are
continuous and strictly positive, and $Z_\rho,V,K_\pm(\rho)$ are bounded.
This proves uniformity, including the limit
$\beta_0=\gamma_0=\sqrt{q/r}$.
\end{proof}

\begin{proposition}
Fix $\rho_{\max}<\infty$. There are compact intervals for state and
control, independent of $T,x_0,c,\rho$, containing every solution (3) for
$0\le\rho\le\rho_{\max}$. Thus the preceding results hold within the
compact-constraint setting of the original question.
\end{proposition}
\begin{proof}
Both exponential factors in (3) are at most one. The coefficient bounds
are uniform on the compact discount interval, so the states and controls
have finite uniform bounds. Choose closed intervals strictly larger than
these bounds. Formula (3) is feasible and the strict-convexity proof applies.
On the chosen state interval, $|\Phi(x)|=|cx|$ is uniformly bounded.
\end{proof}

\section*{Limits of the result}
The two rates are asymmetric: as $\rho\to\infty$,
$\beta_\rho\sim q/(r\rho)$ and $\gamma_\rho\sim\rho$.
Thus even this explicit example does not provide a discount-independent
initial boundary-layer rate on an unbounded discount range.
For the undiscounted static reference $a$, (4) shows a further concrete
failure: if $b\ne0$ and $\epsilon<|\bar x_\rho-a|/2$, then, for all
large $T$, a central interval of length $T-O(1)$ is farther than
$\epsilon$ from $a$. This follows by making each of the two weighted exponential summands in (4)
smaller than $|\bar x_\rho-a|/4$ outside fixed boundary intervals.
The full question remains open here for nonlinear dynamics, active
constraints, general terminal costs, and global existence or selection of
a stationary discounted reference. The derivation proves a substantial
solvable family, not necessary and sufficient conditions for all those cases.

\begin{thebibliography}{9}
\bibitem{fg} T. Faulwasser and L. Gr\"une.
\emph{Turnpike Properties in Optimal Control: An Overview of Discrete-Time and Continuous-Time Results}, arXiv:2011.13670v3, Section 5, p. 25.
\url{https://arxiv.org/pdf/2011.13670v3}.
\end{thebibliography}
\end{document}
